% ============================================================================
% Вывод уравнений ядра: происхождение каждого члена модели и статус обоснования.
% Черновик; цепочка файлов восстановлена целиком: core_model.py (корень) -> daily_core.py ->
% sleep_core2.py -> mass_core6.py -> px/. Пробелы, которых нет ни в одном файле, помечены [вне файлов].
% Компилируется xelatex (как report_ru.tex). Шрифты лежат рядом, в fonts/.
% ============================================================================
\documentclass[11pt,a4paper]{article}
\usepackage{fontspec}
\usepackage{polyglossia}
\setdefaultlanguage{russian}\setotherlanguage{english}
\newcommand{\optserif}{Path=fonts/,Extension=.ttf,UprightFont=*-Regular,BoldFont=*-Bold,ItalicFont=*-Italic,BoldItalicFont=*-BoldItalic,Ligatures=TeX}
\newcommand{\optmono}{Path=fonts/,Extension=.ttf,UprightFont=*-Regular,BoldFont=*-Bold,ItalicFont=*-Italic,BoldItalicFont=*-BoldItalic,Scale=0.9,Ligatures=TeX}
\setmainfont{LiberationSerif}[\optserif]\newfontfamily\cyrillicfont{LiberationSerif}[\optserif]
\setmonofont{LiberationMono}[\optmono]\newfontfamily\cyrillicfonttt{LiberationMono}[\optmono]
\usepackage[a4paper,margin=2.0cm]{geometry}
\usepackage{enumitem,amsmath,amssymb}
\usepackage{booktabs,titlesec,amsthm,array}
\usepackage[hidelinks]{hyperref}
\setlist[itemize]{leftmargin=1.0em,itemsep=0pt,topsep=1pt}
\setlist[enumerate]{leftmargin=1.2em,itemsep=0pt,topsep=1pt}
\setlength{\parskip}{1pt}
\emergencystretch=2em
\titleformat{\section}{\large\bfseries}{\thesection.}{0.4em}{}
\titlespacing*{\section}{0pt}{7pt plus 1pt}{3pt}
\titleformat{\subsection}{\normalsize\bfseries}{\thesubsection.}{0.4em}{}
\titlespacing*{\subsection}{0pt}{5pt}{2pt}
\theoremstyle{remark}
\newtheorem{remark}{Замечание}
\newcommand{\code}[1]{\texttt{\detokenize{#1}}}

\begin{document}

\begin{center}
{\large\bfseries Вывод уравнений ядра}\\[2pt]
{\small Происхождение каждого члена модели и статус его обоснования; черновик}\\[2pt]
{\small 24 октября 2026 г.}
\end{center}
\vspace{-4pt}

\noindent\textbf{Зачем этот документ.} В репозитории нет пошагового вывода уравнений
ядра в одном месте: в \code{mass_core6.py} сказано, что ядро «перенесено из
\code{sleep_core2.py} дословно», и так далее по цепочке. Все звенья цепочки теперь
добавлены в репозиторий, и последнее из них --- \code{core_model.py} --- оказалось
корнем: ни на какой более ранний файл оно не ссылается. Вывод восстановлен из
\code{core_model.py}, брошюры (\code{latex/brochure.tex}) и публикаций. Каждый член
помечен статусом: \emph{измерено}, \emph{выведено из измеренного}, \emph{допущение о
форме} или \emph{вне файлов} (нет ни в одном файле проекта).

\section{Исходная система: измеренная кинетика пероксиредоксина}
\label{sec:selvaggio}
Исходник ядра --- кинетика системы пероксиредоксин--тиоредоксин, измеренная для 15
клеточных линий (Сельваджо, Коэльо, Сальвадор, 2018). Качественно схема реакций такова
(обозначения --- как в брошюре, раздел 8--9):

\begin{enumerate}[label=(Р\arabic*)]
\item окисление: рабочий пероксиредоксин обезвреживает перекись и сам окисляется;
\item починка: окисленная форма за секунды возвращается в рабочее состояние
(тиоредоксиновая система, в избытке);
\item переокисление: если к \emph{уже окисленному} ферменту приходит вторая молекула
перекиси, фермент выключается надолго --- это и есть «второй удар»;
\item медленная починка: переокисленную форму чинит сульфиредоксин, минуты--часы.
\end{enumerate}

Полная модель этой системы даёт два устойчивых состояния лишь у одной клеточной линии
и в узком диапазоне нагрузок; измеренные клетки лежат у границы, но снаружи, и чтобы
попасть внутрь, починку нужно ослабить в $1{,}6$--$22$ раза (брошюра, раздел 9; метка
\emph{допущение или расчёт авторов} --- сам пересчёт упрощённой модели). Из
\emph{структуры} реакций (Р1)--(Р4) следует положительная обратная связь: чем больше
перекиси, тем больше фермента выключено переокислением; чем меньше рабочего фермента,
тем хуже очистка и тем больше перекиси --- при достаточном усилении петля даёт
бистабильность (раздел~\ref{sec:window}).

\section{Редукция к двум переменным и алгебраической перекиси}
\label{sec:reduction}
Из полной системы оставляются две динамические переменные: $x\in[0,1]$ --- доля рабочего
пероксиредоксина, и $c$ --- уровень защиты, поддерживаемый работой NMDA-рецептора.
Перекись $u$ исключается адиабатически (её баланс устанавливается за секунды). Цепочка
переноса: \code{core_model.py} $\to$ \code{daily_core.py} $\to$ \code{sleep_core2.py}
$\to$ \code{mass_core6.py} $\to$ \code{px/}, на каждом шаге дословно; функции
\code{curve}, \code{jac}, \code{window}, \code{states\_at} во всех файлах совпадают,
то есть сама форма ядра не менялась никогда; \code{core_model.py} --- корень: ни на
какой файл оно не ссылается. Ступени редукции:

\begin{itemize}
\item тиоредоксиновая система --- в избытке и не истощается: вместо отдельной переменной
она даёт \emph{эффективную скорость починки} в уравнении для $x$;
\item перекись --- быстрая: вместо дифференциального уравнения --- алгебраический баланс;
\item индукция сульфиредоксина --- медленная (мРНК $\to$ белок, часы--сутки): она и
становится медленной переменной $c$.
\end{itemize}

\subsection{Баланс перекиси}
\begin{equation}\label{eq:balance}
\nu = \varepsilon u + \frac{x\,u}{1+u},\qquad
\varepsilon = e_0\, c^{\gamma}.
\end{equation}
Левая часть --- приток $\nu$ (нагрузка). Первое слагаемое правой части ---
неферментативная очистка (глутатион и прочее): линейная по $u$ с малым коэффициентом.
Вклад защиты в неё параметризован степенной зависимостью $c^{\gamma}$ с разыгрываемым
показателем $\gamma\in[0,1]$ --- \emph{допущение о форме} (простейшая монотонная
параметризация «сколько защиты --- столько и неферментативной очистки»). Второе
слагаемое --- ферментативная очистка рабочим ферментом, с насыщением (константа
Михаэлиса нормирована к единице): это форма реакции (Р1) --- \emph{выведено из
измеренной кинетики} (в исходном \code{core_model.py} строка помечена ровно так:
«выведено из модели Selvaggio и соавт.»).

\subsection{Доля рабочего фермента}
\begin{equation}\label{eq:x}
x' = c^{m}(1-x) - \mu_0\,\frac{x\,u^2}{1+u}.
\end{equation}
Первое слагаемое --- починка окисленной формы (Р2): скорость пропорциональна доле
окисленного фермента $(1-x)$ и усилена защитой. Усиление параметризовано степенью
$c^{m}$ с разыгрываемым показателем $m\in[1,2]$ --- \emph{допущение о форме}; сам факт
«починку делает защита» --- это связь (Р4), но её кинетика в такую форму не измерялась.
Второе слагаемое --- переокисление (Р3): оно требует \emph{второго} удара перекиси по
уже окисленному ферменту, отсюда квадратичность по $u$; знаменатель $1+u$ ---
насыщение --- \emph{выведено из измеренной кинетики} (в \code{core_model.py} помечено
«выведено», время --- в единицах времени восстановления фермента). Множитель $\mu_0$ ---
масштаб скорости переокисления относительно починки, разыгрывается в $[0{,}01;10]$.

\subsection{Уровень защиты: допущение о форме}
\begin{equation}\label{eq:c}
c' = \frac{1}{T}\Bigl(N + \frac{\beta\,u}{K_B+u} - c\Bigr),\qquad
N = \frac{N_0}{1+u/K_N}.
\end{equation}
Это уравнение \emph{целиком --- допущение о форме}, и это единственный член ядра, не
опирающийся на измеренную кинетику. Обосновано в нём только следующее:

\begin{itemize}
\item связь «рецептор поддерживает защиту» --- из статей (синаптическая NMDA-активность
включает гены антиоксидантной защиты; Пападия и соавт., 2008). В модели её несёт
слагаемое $N$;
\item подавление этой части защиты перекисью --- $N_0/(1+u/K_N)$: монотонное насыщающееся
подавление; $N_0\in[0{,}3;1]$ изображает врождённо слабый рецептор;
\item индукция защиты самой перекисью --- $\beta u/(K_B+u)$: монотонное насыщающееся
включение (путь Nrf2);
\item релаксация первого порядка со временем $T$ --- время жизни защитных белков.
\end{itemize}

Почему выбрана именно такая форма: это простейшие монотонные насыщающиеся функции,
дающие два конкурирующих вклада перекиси в защиту (подавление через $N$ и индукцию
через $\beta$) и ограниченную $c$; бистабильность при этом делает петля $x$--$u$
(раздел~\ref{sec:window}), а $c$ лишь сдвигает границы окна. Время в модели измеряется
в единицах времени починки фермента, поэтому $T$ --- это отношение времени жизни
защитных белков к времени починки. В исходном \code{core_model.py} уравнение помечено
как «защита следует за активностью рецептора», $N$ --- «форма --- допущение», и дана
семантика всех девяти параметров: $e_0$, $\mu_0$ --- быстрый слой; $m$, $\gamma$ --- как
восстановление и обходная очистка зависят от защиты; $K_N$, $N_0$ --- рецептор;
$\beta$, $K_B$ --- индукция защиты стрессом; $T$ --- отношение масштабов времени.
Числовые \emph{границы} диапазонов ($e_0\in[0{,}005;0{,}1]$ и т.\,д.) не обоснованы
ни в одном файле [вне файлов]; известно лишь, что $T$ в \code{core_model.py} был
$[10;300]$ и сужен до $[20;80]$ в \code{daily_core.py} ради шага счёта.

\begin{remark}
Единственная задокументированная правка выборки: $\beta$ задаётся не сама по себе, а
кратностью $\beta=N_0\cdot U(0,3)$, потому что измеренная индукция белка сульфиредоксина
достигает $2{,}85$ раза, а прежняя выборка $\beta\in[0,1]$ при $N_0\in[0{,}3;1]$ такой
кратности почти не допускала (\code{px/core.py}, \code{mass_core6.py}, раздел 4е).
\end{remark}

\section{Откуда берётся окно: бистабильность ядра}
\label{sec:window}
Бистабильность --- следствие петель (Р1)--(Р3), и её видно прямо из
уравнений~\eqref{eq:balance}--\eqref{eq:x}. При замороженной защите $c$ стационарная
доля рабочего фермента есть
\begin{equation}
x^*(u) = \frac{c^{m}}{c^{m} + \mu_0\, u^2/(1+u)},
\end{equation}
--- монотонно убывающая функция перекиси. Подстановка в
баланс~\eqref{eq:balance} даёт требуемый приток как функцию $u$:
\begin{equation}\label{eq:null}
\nu(u) = \varepsilon u + \frac{c^{m}}{c^{m} + \mu_0\,u^2/(1+u)}\cdot\frac{u}{1+u}.
\end{equation}
Пока неферментативная очистка мала ($\varepsilon\ll1$), функция $\nu(u)$
\emph{немонотонна}: при малых $u$ очистка почти ферментативная и растёт с $u$, при
больших --- фермент выбит переокислением, и $\nu(u)$ растёт снова; между ними ---
убывающий участок. Немонотонность $\nu(u)$ и есть пара седло-узловых складок: в
интервале нагрузок между локальным максимумом и минимумом~\eqref{eq:null} существуют
два устойчивых стационара («защита работает», «защита выключена») --- окно гистерезиса.

\emph{Точный критерий быстрого слоя (задокументирован в \code{core_model.py}).} При
замороженной защите, после перенормировки $x^*(u)=1/(1+\mu\,u^2/(1+u))$ с
$\mu=\mu_0/c^{m}$, баланс даёт $\nu(u)=\varepsilon u+u/(1+u+\mu u^2)$, откуда
\begin{equation}
\frac{d\nu}{du}=\varepsilon+\frac{1-\mu u^2}{(1+u+\mu u^2)^2},
\end{equation}
и у быстрого слоя есть два состояния тогда и только тогда, когда
\begin{equation}\label{eq:phi}
\varepsilon<\Phi(\mu),\qquad
\Phi(\mu)=\max_{u>0}\frac{\mu u^2-1}{(1+u+\mu u^2)^2}
\end{equation}
--- это функция \code{Phi} из \code{core_model.py}; она же и есть «упрощение до одного
уравнения, дающее точный критерий двух состояний» из брошюры. Применение критерия к 15
измеренным клеточным линиям, давшее «починку нужно ослабить в $1{,}6$--$22$ раза»,
отдельным файлом не сохранено [вне файлов].

\emph{Классы поведения и оценки из \code{core_model.py}.} Система на плоскости
ограничена, поэтому при нагрузке, где нет ни одного устойчивого состояния, она обязана
выйти на цикл --- класс «колебания» определён так и сверен прямым интегрированием.
Там же записаны и проверены две аналитические оценки: порог нагрузки падает как
$N_0^{\,m/2}$ (правило П1) и замедление возврата у границы окна идёт по закону
седло-узлового призрака с показателем $-1/2$ (правило П4).
Медленная защита $c$ эти границы только двигает, не создавая и не убирая окно. Реализованный
в коде критерий --- ровно этот: \code{window} ищет убывающий участок $\nu(u)$
(\code{diff(nu) < 0}), а \code{states\_at} отбирает корни с проверкой следа и определителя
якобиана. Набор принимается, только если окно есть и в его среднегеометрической середине
два устойчивых состояния --- это «потенциальные больные» модели. Проверочный набор с
известным окном существует с самого начала (\code{sleep_core2.py}, раздел 0):
$e_0=0{,}02$, $\mu_0=0{,}3$, $m=2$, $\gamma=0{,}5$, $K_N=3$, $N_0=0{,}9$,
$\beta=0{,}1$, $K_B=10$, $T=50$ --- на нём же шаг сверялся с точным решателем.

\section{Что добавлено поверх ядра, и откуда каждая форма}
\label{sec:extensions}

\subsection{Две ступени защиты}
\begin{equation}
\dot m=\frac{1}{\tau_m}\Bigl(\beta\frac{u}{K_B+u}-m\Bigr),\qquad
\dot s=\frac{1}{\tau_s}(m-s).
\end{equation}
\emph{Форма} --- два звена первого порядка подряд (мРНК $\to$ белок), допущение,
мотивированное биологией. \emph{Постоянные} --- из сырых измерений пути Nrf2 (876 точек
по сульфиредоксину, 17 условий, 44 часа): $\tau_m=12$\,ч --- по мРНК (8,1--9,6\,ч в трёх
обработках), $\tau_s=63$\,ч --- по белку; спад белка в опубликованных работах не
измерен, по разным способам обработки выходит 24--88\,ч, рабочее значение --- середина,
и это допущение, а не измерение. Спад считался по разности «доза минус контроль»,
потому что контроль сам даёт горб и без вычитания завышает спад
(\code{mass_core6.py}, раздел 4е; брошюра, «Медленная защита»).

\subsection{Прайминг}
\begin{equation}
\dot p = \Bigl(\frac{u}{K_B+u}-p\Bigr)\Big/\tau_p,\qquad
\beta_{\mathrm{eff}}=\beta\,(1+A\,p),\qquad A=2.
\end{equation}
Потолок усиления $\times3$ измерен (повторное воздействие через 24\,ч даёт наработку
фермента до трёх раз выше; Бишофф и соавт.). Постоянная памяти $\tau_p$ не измерена ---
в опыте проверен один интервал, 24\,ч, --- поэтому перебирается. Форма следа ---
допущение (\code{px/rules/f19.py}, брошюра, «Прайминг»).

\subsection{Асимметрия подъёма и спада}
Постоянные подъёма и спада второй ступени разведены: $\tau_{\uparrow}\ne\tau_{\downarrow}$,
$r=\tau_{\downarrow}/\tau_{\uparrow}$. Полный вывод следствия --- в брошюре,
«Асимметричная защита»: для двухуровневого суточного входа, высокого долю $f$ суток,
стационарное среднее утекающего интегратора стоит в точке
\begin{equation}
w(r)=\frac{f\,r}{f\,r+1-f},
\end{equation}
где $f=0{,}705$ --- измеренная доля бодрствования (не назначенная), поэтому в
предсказании нет подогнанных чисел. Равенство постоянных было невыведенным допущением,
которое версия 5 сняла.

\subsection{Петля ёмкости}
Пока клетка повреждена, её $N_0\to N_0(1-a_n)$ (потеря 1\,сут, быстрая --- не
перебирается), возврат медленный ($\tau_{\uparrow}=15\dots240$\,сут, перебирается).
Обоснование биологическое, а не кинетическое: окислительное повреждение PV-клеток
снижает экспрессию самого парвальбумина, рвёт перинейрональную сеть, бьёт по
митохондриям --- клетка теряет \emph{ёмкость}, а не получает больше нагрузки.
Асимметрия быстрой потери и медленного возврата заимствована из результата Ф23
(\code{px/rules/f28.py}).

\subsection{Дофаминовая петля}
Один скаляр $d$ (тонус), движимый растормаживанием $E/E_r-1$ или долей $D$, с
постоянной $\tau_d=7$\,сут, входит в нагрузку множителем $(1+a_d\,d)$. Здесь честно
записано, что вывода нет: «литература подтверждает направление --- катаболизм дофамина
даёт перекись, --- но не даёт ни одного количественного значения: ни скорости
образования, ни постоянной времени петли». Поэтому правило Ф26 задаёт структурный, а
не параметрический вопрос (\code{px/rules/f26.py}, комментарии в \code{px/model.py}).

\subsection{Историческая форма нагрузки и сна}
\label{sec:history}
До переключателя Филлипса--Робинсона сон был двухпороговым: давление сна растёт к 1
с постоянной 18{,}2\,ч и спадает к 0 с постоянной 4{,}2\,ч (Даан, Беерсма, Борбели,
1984 --- уже там); бодрствование кончается при $S=H$, сон --- при $S=L$; пороги
$H=0{,}67$ и $L=0{,}17$ с суточной модуляцией $0{,}12$ (эта схема --- из
\code{daily_core.py} и \code{sleep_core2.py}). Связь «повреждение мешает заснуть» ---
$H_{\mathrm{эфф}}=H+k\,d$, с потолком $0{,}98$.
Нагрузка имела вид
\begin{equation}
\nu(t)=\nu_{\mathrm{ср}}\,\frac{1+A\,S(t)}{1+A\,\bar S_0}\,\exp(\eta(t)),\qquad A=0{,}2,
\end{equation}
где $\bar S_0$ --- среднее давление сна за свободный цикл при $k=0$. Отсюда родом обе
нынешние величины: суточный размах нагрузки \code{DSWING}$=0{,}20$ --- это тот самый
$A=0{,}2$, перенесённый с давления сна на сомногены в \code{mass_core6.py}; а член
$-k_d D$ в притоке к сонному ядру --- прямой наследник $H_{\mathrm{эфф}}=H+k\,d$.
\emph{Форма} нагрузки --- произведение суточного множителя и логнормального шума --- была
принята тогда же; вывод её, впрочем, задокументирован частично --- см. ниже.

\emph{Вывод амплитуды $A=0{,}2$.} В \code{daily_core.py} записаны три литературных
опоры: захват глюкозы корой мыши во сне ниже на 16\,\%; потребление кислорода у
человека при засыпании ниже на 21\,\%; метка окисления ДНК в PV-клетках крысы после
6\,ч нарушенного сна выше на 25\,\%. Три линии сходятся на отношении $\approx1{,}2$ ---
отсюда рабочий размах $A=0{,}2$ (ныне \code{DSWING}$=0{,}20$). Честная оговорка там же:
переокисленный пероксиредоксин --- главная переменная модели --- по времени суток в коре
никем не измерен, поэтому $A$ опирается только на косвенные величины.

\emph{Вывод периода суток $r$.} Сутки в единицах модели --- это $r\,T$, $r\in[2,4]$;
основание записано в \code{daily_core.py}: единица времени --- время починки фермента,
$T$ --- время жизни защитных белков в этих единицах, а $24$\,ч / (время жизни белков
6--12\,ч) $=2$--$4$. История диапазона $T$ тоже записана: в \code{core_model.py} он был
$[10;300]$, а в \code{daily_core.py} сужен до $[20;80]$ --- при большом $T$ сутки
длинны, и шаг счёта пришлось бы дробить без пользы.

\emph{Медленный шум.} Время корреляции шума $\tau_c=0{,}3\,T$ выбрано по результату
прошлого расчёта: быстрый шум правдоподобной силы оказался слишком слаб, чтобы
перебрасывать систему (\code{daily_core.py}).

\emph{Порог эпизода.} Определение «повреждённая доля $d=1-x$ выше середины между двумя
устойчивыми ветвями, взятой в середине окна» --- то есть нынешний порог
\code{thr}$=0{,}5(x_{\mathrm{lo}}+x_{\mathrm{hi}})$ --- впервые записано в
\code{daily_core.py}.

Волна процесса S строилась по постоянным Даана (18{,}2/4{,}2\,ч), нормировалась к
размаху $0$--$1$, а среднее держалось постоянным делением на $(1+A\,\bar S)$;
самопроверка снимала постоянные времени обратно из построенной волны по её собственным
асимптотам.

\subsection{Слои сна и коры}
Переключатель Филлипса--Робинсона (набор «new nominal»; Скелдон и соавт., 2014) и пара
Вильсона--Коуэна с $g_{\mathrm{dep}}=0{,}4$ (Вольман и соавт., 2011) взяты из статей
целиком; здесь они не выводятся --- их сверка с опубликованными числами приведена в
брошюре (раздел «Переключатель сна») и в обоих отчётах.

\section{Сводка статусов}
\label{sec:table}
\begin{center}\footnotesize
\begin{tabular}{@{}>{\raggedright\arraybackslash}p{0.38\textwidth}>{\raggedright\arraybackslash}p{0.20\textwidth}>{\raggedright\arraybackslash}p{0.36\textwidth}@{}}
\toprule
Член / параметр & Статус & Где обосновано\\
\midrule
баланс $\nu=\varepsilon u+xu/(1+u)$, форма $xu/(1+u)$ & выведено из измеренного &
кинетика (Р1); пометка «выведено» --- в \code{core_model.py}\\
$\varepsilon=e_0c^{\gamma}$, $\gamma\in[0,1]$ & допущение о форме & степенная
параметризация вклада защиты\\
$x$-уравнение: починка $c^{m}(1-x)$ & форма --- допущение; связь --- из (Р4) & $m\in[1,2]$
разыгрывается\\
$x$-уравнение: переокисление $\mu_0 xu^2/(1+u)$ & выведено из измеренного & второй удар
(Р3), насыщение; пометка «выведено» --- в \code{core_model.py}\\
$c$-уравнение целиком & допущение о форме & связь «рецептор $\to$ защита»: Пападия и
соавт., 2008; пометка «форма --- допущение» --- в \code{core_model.py}\\
$N_0\in[0{,}3;1]$ & параметр, смысл из гипотезы & врождённо слабый рецептор\\
$\beta=N_0\cdot U(0,3)$ (прежде $\beta\in[0,1]$) & измерено & индукция до $2{,}85\times$; правка --- в \code{mass_core6.py}, 4е\\
диапазоны $e_0$, $\mu_0$, $m$, $\gamma$, $K_N$, $K_B$ & семантика задана, границы ---
без обоснования & \code{core_model.py}; числовые границы не обоснованы нигде
[вне файлов]; $T$: было $[10;300]$, сужено в \code{daily_core.py}\\
критерий $\varepsilon<\Phi(\mu)$ & выведено & \code{core_model.py}, функция \code{Phi}\\
оценки $N_0^{\,m/2}$ и призрака $-1/2$ & выведено и проверено & правила П1 и П4,
\code{core_model.py}\\
суточный размах $A=0{,}2$ (\code{DSWING}) & выведено из трёх измерений &
16\,\%/21\,\%/25\,\% $\to$ $\approx1{,}2$, \code{daily_core.py}\\
период суток $r\,T$, $r\in[2,4]$ & выведено & 24\,ч/(6--12\,ч) $=2$--$4$,
\code{daily_core.py}\\
порог эпизода $\mathrm{thr}=0{,}5(x_{\mathrm{lo}}+x_{\mathrm{hi}})$ & определение &
середина между ветвями в середине окна, \code{daily_core.py}\\
$\tau_m=12$\,ч, $\tau_s=63$\,ч & сырые данные; 63\,ч --- допущение & \code{mass_core6.py},
раздел 4е\\
прайминг: $A=2$ & измерено & Бишофф и соавт.; $\tau_p$ не измерена, перебирается\\
асимметрия: $w(r)$ & выведено & брошюра, «Асимметричная защита»\\
ёмкость: $a_n$, $\tau_{\uparrow}$, $\tau_{\downarrow}$ & допущение, биологическое
обоснование & \code{f28.py}; асимметрия --- из Ф23\\
дофамин: $a_d$, $\tau_d$ & не измерены по литературе & \code{f26.py}, \code{model.py}\\
сон, кора & взяты из статей целиком & Скелдон 2014; Вольман 2011\\
\bottomrule
\end{tabular}
\end{center}

\section{Что осталось закрыть}
\label{sec:gaps}
Цепочка восстановлена целиком: \code{core_model.py} --- корень (ни на какой файл не
ссылается), далее \code{daily_core.py} $\to$ \code{sleep_core2.py} $\to$
\code{mass_core6.py} $\to$ \code{px/}, на каждом шаге дословно. Добавленные файлы
закрыли: (i) одноклеточные числа $0{,}009$--$0{,}013$ (\code{diag_gradual.py} ---
диагностика); (ii) вывод $A=0{,}2$, $r\in[2,4]$, историю $T$, выбор медленного шума и
порог эпизода (\code{daily_core.py}); (iii) точный критерий $\Phi(\mu)$, оценки
$N_0^{\,m/2}$ и призрака $-1/2$, класс «колебания», семантику всех девяти параметров и
пометки «выведено»/«форма --- допущение» (\code{core_model.py}).

\emph{Ниже --- три оставшихся пробела, каждый с пометкой: что известно из файлов, чего
нет, и что является реконструкцией, а не источником.}

\subsection{Пробел 1: числовые границы диапазонов}
\label{sec:gap1}
Что известно: в \code{core_model.py} задана семантика каждого из девяти параметров
(раздел~\ref{sec:reduction}), а история диапазона $T$ задокументирована
($[10;300]\to[20;80]$). Сами границы
\begin{equation*}
\begin{gathered}
e_0\in[0{,}005;0{,}1],\quad \mu_0\in[0{,}01;10],\quad m\in[1,2],\quad
\gamma\in[0,1],\\
K_N\in[0{,}1;100],\quad N_0\in[0{,}3;1],\quad K_B\in[1;100]
\end{gathered}
\end{equation*}
заданы в \code{RANGES} с самого первого файла без построчного обоснования.
\emph{[пробел --- вне файлов.]} Ниже --- только согласованность границ со структурой
уравнений, выведенная здесь, а не взятая из файлов \emph{[реконструкция]}: $e_0$ мал по
построению (бистабильность требует малой неферментативной очистки,
раздел~\ref{sec:window}); $\mu_0$ --- главный рычаг критерия~\eqref{eq:phi}, и
диапазон $[0{,}01;10]$ заведомо накрывает переход; $m\in[1,2]$ --- усиление починки
защитой от линейного до квадратичного; $\gamma\in[0,1]$ --- обходная очистка от
независимой от защиты до пропорциональной ей; $K_N$, $K_B$ --- положения насыщения
подавления и индукции, взяты широко, «с запасом»; $N_0=1$ --- здоровый рецептор,
нижняя граница $0{,}3$ --- ослабление втрое. Откуда взяты именно эти числа --- нигде не
записано.

\subsection{Пробел 2: шаги редукции из полной модели Сельваджо}
\label{sec:gap2}
Что известно: пометки «выведено из модели Selvaggio и соавт.» (баланс) и «выведено»
(уравнение для $x$) в \code{core_model.py}; качественная схема реакций (Р1)--(Р4) и
три ступени редукции --- в разделе~\ref{sec:reduction} этого документа.
\emph{[пробел --- вне файлов.]} Нигде не записано построчное соответствие между
скоростями полной модели (окисление, починка тиоредоксином, переокисление вторым
ударом, починка сульфиредоксином) и членами ядра: какие комбинации констант свернулись
в $e_0$, $\mu_0$, $m$, $\gamma$ и как из двухударной схемы получилась именно форма
$\mu_0\,x\,u^2/(1+u)$. Реконструкция формы дана в разделе~\ref{sec:reduction}
\emph{[реконструкция]}, но это вывод из структуры, а не из авторских заметок.

\subsection{Пробел 3: применение $\Phi(\mu)$ к измеренным линиям}
\label{sec:gap3}
Что известно: критерий~\eqref{eq:phi} из \code{core_model.py} и утверждение брошюры
(раздел 9): полная модель даёт два состояния лишь у одной из 15 клеточных линий и в
узком диапазоне нагрузок; измеренные клетки лежат у границы, но снаружи, и чтобы
попасть внутрь, починку нужно ослабить в $1{,}6$--$22$ раза.
\emph{[пробел --- вне файлов.]} Самого расчёта нет: ни параметров 15 линий, ни способа
«ослабить починку» в терминах $\mu$, ни полученных множителей с их разбросом.
\emph{[реконструкция]} из~\eqref{eq:phi}: ослабить починку --- значит увеличить
$\mu=\mu_0/c^{m}$; «$1{,}6$--$22$ раза» --- это, по-видимому, разброс требуемого
увеличения $\mu$ по линиям, при котором начинает выполняться
$\varepsilon<\Phi(\mu)$; но файла с этим расчётом нет, и числа остаются
невоспроизводимыми из репозитория.

Эти три пробела не влияют на работу модели --- они касаются только происхождения
чисел; при необходимости их можно закрыть отдельным расчётом.

\smallskip
\begin{thebibliography}{9}\small
\itemsep0pt
\bibitem{selvaggio} G. Selvaggio, P. Coelho, A. Salvador. The peroxiredoxin--thioredoxin
system: kinetic modeling. \emph{Redox Biology}, 2018.
\bibitem{papadia} S. Papadia et al. Synaptic NMDA receptor activity boosts intrinsic
antioxidant defenses. \emph{Nat. Neurosci.} 11, 476 (2008).
\bibitem{bischoff} Bischoff et al. Sulfiredoxin induction kinetics (raw data), \emph{Sci. Rep.};
priming ceiling $\times3$.
\bibitem{skeldon} A. C. Skeldon, D.-J. Dijk, G. Derks. \emph{PLoS One} 9, e103877 (2014).
\bibitem{volman} V. Volman, M. M. Behrens, T. J. Sejnowski. \emph{J. Neurosci.} (2011).
\bibitem{daan} S. Daan, D. G. M. Beersma, A. A. Borb\'ely. \emph{Am. J. Physiol.} 246, R161 (1984).
\end{thebibliography}

\end{document}
